%----------------------------------------------------------------------------------------
% NOTATIE EN FIGUREN VAN DE CURSUS GRAFEN (2de graad)
% Staat in de preambule en niet in een hoofdstuk: ook de lijsten met oplossingen gebruiken dit.
% Vertrokken van notatie.tex van Integralen; wat hier niet nodig is, is geschrapt.
%----------------------------------------------------------------------------------------
% Bewaarde grafiek, geplaatst door maak-pdf.py: bestand en pagina in figuren-cache/
\newcommand{\grafiekbestand}[2]{\includegraphics[page=#2,trim=3bp 3bp 3bp 3bp]{figuren-cache/#1}}

% De graad van een knoop: \gr{A}
\DeclareMathOperator{\grop}{gr}
\newcommand{\gr}[1]{\grop(#1)}

% Zoals stappen, met eigen kolombreedtes: \begin{stappenvrij}{6.2cm}{3.8cm}. Voor naast een figuur.
\newenvironment{stappenvrij}[2]{%
  \par\renewcommand{\arraystretch}{1.25}%
  \noindent\begin{tabular}{@{}>{$\displaystyle}p{#1}<{$}@{\hspace{3mm}}>{\RaggedRight\small\itshape\arraybackslash}p{#2}@{}}%
}{\end{tabular}\par}

% Tekst of berekening links, figuur rechts: \naast{links}{rechts}
\newsavebox{\naastbox}
\newcommand{\naast}[3][]{%
  \par\noindent\sbox{\naastbox}{#3}%
  \begin{minipage}[t]{\dimexpr\linewidth-\wd\naastbox-3mm\relax}\vspace{0pt}#2\end{minipage}\hfill
  \begin{minipage}[t]{\wd\naastbox}\vspace{0pt}\usebox{\naastbox}\end{minipage}\par\addvspace{1mm}}
% Een rij figuren, gecentreerd
\newcommand{\figuren}[1]{\par\addvspace{1mm}{\centering #1\par}\addvspace{1.5mm}}
% Opgave met een figuur rechts: \begin{metgrafiek}{tikz} opgave en deelvragen \end{metgrafiek}
\newsavebox{\grafiekbox}
\newenvironment{metgrafiek}[2][]{%
  \sbox{\grafiekbox}{#2}%
  \begin{minipage}[t]{\dimexpr\linewidth-\wd\grafiekbox-3mm\relax}%
}{%
  \end{minipage}\hfill
  \begin{minipage}[t]{\wd\grafiekbox}\vspace{-2mm}\usebox{\grafiekbox}\end{minipage}\par}

% GRAFEN TEKENEN, binnen een tikzpicture.
% \knopen{A/0/0, B/2/0, C/1/1.5}   knopen met hun naam erin, op de gegeven plaats (x/y in cm, of de schaal van de figuur)
% \bogen{A/B, B/C}                 bogen tussen knopen (een rechte lijn); een kromme boog: \draw[boog] (A) to[bend left=25] (B);
% \pijlen{A/B, C/A}                gerichte bogen, van de eerste naar de tweede knoop
% \gbogen{A/B/5, B/C/3}            gewogen bogen, met het gewicht op een wit vlakje in het midden
% \gpijlen{A/B/5}                  gewogen gerichte bogen
% \knoopgraad[hoek]{A}{3}          een getal naast een knoop (de graad), in de kleur van het hoofdstuk; hoek in graden (standaard 45)
% Stijlen: knoop (cirkel met naam), punt (kleine stip zonder naam), boog, pijl, gewicht, accentboog (dik, voor een pad).
\tikzset{
  knoop/.style={circle,draw=accent,thick,fill=white,inner sep=0pt,minimum size=5.4mm,font=\footnotesize,text=black},
  punt/.style={circle,fill=accent,inner sep=0pt,minimum size=4.5pt},
  boog/.style={black!65,thick},
  pijl/.style={boog,-{Stealth[length=2.2mm,width=1.8mm]},shorten >=.5pt},
  gewicht/.style={fill=white,inner sep=1pt,font=\scriptsize,text=black},
  accentboog/.style={accent,line width=2.2pt},
  graadlabel/.style={text=accent,font=\scriptsize\bfseries,inner sep=1.5pt},
}
\newcommand{\knopen}[1]{\foreach \kn/\kx/\ky in {#1}{\node[knoop] (\kn) at (\kx,\ky) {$\kn$};}}
\newcommand{\bogen}[1]{\foreach \ka/\kb in {#1}{\draw[boog] (\ka) -- (\kb);}}
\newcommand{\pijlen}[1]{\foreach \ka/\kb in {#1}{\draw[pijl] (\ka) -- (\kb);}}
\newcommand{\gbogen}[1]{\foreach \ka/\kb/\kw in {#1}{\draw[boog] (\ka) -- node[gewicht] {$\kw$} (\kb);}}
\newcommand{\gpijlen}[1]{\foreach \ka/\kb/\kw in {#1}{\draw[pijl] (\ka) -- node[gewicht] {$\kw$} (\kb);}}
\newcommand{\knoopgraad}[3][45]{\node[graadlabel] at ($(#2)+(#1:5mm)$) {#3};}
